\chapter{Software und Softwarearchitektur}
\label{ch:01_software-and-software-architecture}

\begin{universityTask}
    \begin{enumerate}[label=(\alph*)]
        \item Beschreiben Sie, was man unter \q{Software} versteht. Gehen Sie dabei zunächst auf die Unterschiede zwischen Software im engeren Sinne und Software im weiteren Sinne ein. Im Anschluss beschreiben Sie bitte Unterschiede zwischen Algorithmen und Programmen.
            \universityPoints{8}
        \item Bearbeiten Sie die nachfolgende Fallstudie zur Komplexität von Softwaresystemen. Analysieren Sie zunächst die beteiligten Systemelemente der Fallstudie. Bilden Sie dazu die maximale Anzahl an Systembeziehungen in Form eines sogenannten Point-to-Point-Szenarios ab. Berechnen Sie darauf aufbauend die Gesamtkomplexität des Gesamtsystems nach der O-Notation und stellen Sie ihre Berechnung grafisch dar. \linebreak
            \begin{quotation}
                \begin{itshape}
                    Ein Kunde bestellt zum ersten Mal im Online-Shop Wasserschlauch.de. Im Rahmen des Bestellprozesses wählt der Nutzer die Möglichkeit, sich mittels eines ihm bekannten und häufig genutzten \q{Login-Verfahrens} eines sozialen Netzwerkes (beispielsweise mit Facebook, Twitter oder XING) zu authentifizieren. In einem weiteren Schritt werden die vom Kunden ergänzten und die aus dem sozialen Netzwerk übernommenen Kundendaten mit der Kundendatenbank von Wasserschlauch.de abgeglichen. Zudem erfolgt eine SCHUFA-Anfrage, die die Bonität des neuen Kunden absichert. Um dem Kunden ein voraussichtliches Lieferdatum für seine Bestellung anzeigen zu können, werden im nächsten Schritt das Warenwirtschaftssystem von Wasserschlauch.de wie auch das Dispositionssystem des Logistikdienstleisters abgefragt und dem Nutzer im Online-Shop als kumuliertes Ergebnis angezeigt. Der Kunde kann zudem entscheiden, ob er Meilen oder Punkte im Rahmen des Kaufprozesses sammeln möchte -- entscheidet er sich hierfür und trägt er seine Kartennummer ein, werden die Informationen zu den gesammelten Meilen oder Punkten an das Kundenverwaltungssystem des Loyalitätsdienstleisters übermittelt. Zuletzt erfolgt unter Einsatz einer finanzwirtschaftlichen Anwendung die Erstellung der Rechnung an den Kunden.
                \end{itshape}
            \end{quotation}
            \universityPoints{12}
            \pagebreak
        \item Benennen Sie die zentralen Herausforderungen, denen sich moderne Formen der Softwareentwicklung stellen müssen, und formulieren Sie Möglichkeiten zur Lösung.
            \universityPoints{5}
        \item Was versteht man unter Objekten, Objektklassen, Objektattributen und -operationen? Definieren und beschreiben Sie die Begriffe kurz und grenzen Sie sie voneinander ab.
            \universityPoints{8}
        \item Definieren Sie den Begriff "Softwarearchitektur" und beschreiben Sie die Unterschiede zwischen fachlicher, technischer und Plattform-Architektur.
            \universityPoints{7}
    \end{enumerate}
\end{universityTask}

\section{Software}
\label{sec:01-01_software}

Der Text enthält die Lösung der obigen Aufgabenstellung.

\section{Komplexität von Softwaresystemen}
\label{sec:01-02_complexity-of-software-systems}

Der Text enthält die Lösung der obigen Aufgabenstellung.

\section{Herausforderungen der Softwareentwicklung}
\label{sec:01-03_challenges-of-software-development}

Der Text enthält die Lösung der obigen Aufgabenstellung.

\section{Objekte, Objektklassen, Objektattribute und -operationen}
\label{sec:01-04_objects-classes-attributes-and-operations}

Der Text enthält die Lösung der obigen Aufgabenstellung.

\section{Softwarearchitektur}
\label{sec:01-05_software-architecture}

Der Text enthält die Lösung der obigen Aufgabenstellung.